% !PRINTED 19
\weberpage{19}

Die einzige Zahl $0$ hat den absoluten Werth $0$. Jede positive Zahl
kommt bei unendlich vielen complexen Zahlen als absoluter Werth vor und
die Bildpunkte aller Zahlen mit demselben absoluten Werthe liegen auf
einem Kreise, dessen Mittelpunkt im Coordinatenanfangspunkte liegt.

Zwei imaginäre Zahlen, die sich nur durch das Vorzeichen von $i$
unterscheiden, also $x+yi$ und $x-yi$, heissen conjugirt imaginär. Ihr
Product ist das Quadrat des absoluten Werthes von jeder von ihnen.

Wenn von zwei conjugirten imaginären Zahlen die eine gleich Null ist,
so ist auch die andere gleich Null, und man kann also in jeder
richtigen Zahlengleichung $i$ durch $-i$ ersetzen, ohne dass die
Richtigkeit gestört wird.

Wir wollen noch den oft angewandten Satz anführen, dass der absolute
Werth einer Summe zweier von Null verschiedener complexer Zahlen
niemals grösser ist, als die Summe der absoluten Werthe der Summanden,
und nur dann gleich, wenn das Verhältniss (der Quotient) beider
Summanden reell und positiv ist. Sei nämlich
\[
 z=x+yi,
 \qquad c=a+bi,
 \qquad Z=(x+a)+(y+b)i,
\]
\[
 \varrho^2=x^2+y^2,
 \qquad r^2=a^2+b^2,
 \qquad R^2=(x+a)^2+(y+b)^2,
\]
dann ist
\[
 (r+\varrho-R)(r+\varrho+R)
   =(r+\varrho)^2-R^2=2(r\varrho-ax-by);
\]
das ist sicher positiv, wenn $ax+by\leqq 0$ ist. Wenn aber $ax+by>0$
ist, so folgt aus
\[
 r^2\varrho^2-(ax+by)^2=(ay-bx)^2,
\]
dass $r\varrho\geqq ax+by$ ist, und nur dann gleich, wenn $ay-bx=0$.
Daraus also ergiebt sich, dass $r+\varrho>R$ ist und nur in dem
besonderen Falle $r+\varrho=R$, wenn $ay-bx=0$, $ax+by>0$, woraus das
Gesagte folgt.

Bei der geometrischen Darstellung ist dieser Satz ein Ausdruck dafür,
dass in einem Dreieck eine Seite kleiner ist, als die Summe der beiden
anderen. Der Satz hat noch die andere Folge, dass der absolute Werth
einer Summe nicht kleiner sein kann, als die Differenz der absoluten
Werthe der Summanden. Denn wenden wir den vorigen Satz auf die Summe
$c=Z-z$ an, so folgt $r\leqq R+\varrho$ oder
\[
 R\geqq r-\varrho.
\]
%CONT>>
